
% SGA 3，第 VII 篇_B；英文重建本第一轮的中文译本。
% 范围：依据文献局部页 55--56；第 2.5 节至命题及注 2.5.A，止于 2.5.1 前。
% 合并阅读本页 568--569；重版印刷页 557--558；原版边码 522。
% 编者注 95--96。依据：Polo--Gille Exp7B-13oct24.pdf。
% SHA-256 D303690E0195654BBE7C3DAA5D89EA23E24CD709A3575217234AC19AE8836694。
% 参照了 Jacob Reinhold 的 CC BY 4.0 Markdown，作为对照来源。本段无图。
\subsection*{2.5}
\phantomsection
\label{sga3:VIIB:2.5}
\addcontentsline{toc}{subsection}{2.5. 平展形式群}
若形式 \(k\)-群的底形式簇是平展的，就称其为\emph{平展}形式群\textup{（参见\ \SGARefLink{sga3:VIIB:1.6}{1.6}）}。这种形式群的结构非常简单。设 \(k\) 是局部环，\(\kappa\) 是其剩余域，\(\kappa_s\) 是 \(\kappa\) 的可分闭包，并令
\[
  \Gamma=\operatorname{Gal}(\kappa_s/\kappa)
\]
为拓扑 Galois 群。所谓 \emph{\(\Gamma\)-集合}，是指带有 \(\Gamma\) 的连续作用的集合 \(E\)；等价地，每个元素 \(x\in E\) 的稳定子都是 \(\Gamma\) 的开子群。
对每个形式 \(k\)-簇 \(X\)，置
\[
  X(\kappa_s)
  =
  \varinjlim_{\ell}X(\ell),
\]
其中 \(\ell\) 遍历 \(\kappa_s\) 中包含的 \(\kappa\) 的有限扩张。\footnote{%
\label{sga3:VIIB:nde:95}若 \([\kappa_s:\kappa]=\infty\)，则 \(\kappa_s\) 不是预有限 \(\kappa\)-代数，所以上述记号是滥用。以下细节也是编者补充的。}
群 \(\Gamma\) 连续作用于每个 \(X(\ell)\)，从而也作用于 \(X(\kappa_s)\)。此外，若
\[
  X_\kappa=X\widehat\otimes_k\kappa
  \qquad\text{\textup{(参见\ 1.6.C)}},
\]
则对每个 \(\ell\) 有 \(X(\ell)=X_\kappa(\ell)\)，故
\[
  X(\kappa_s)=X_\kappa(\kappa_s).
\]
现在设 \(X\) 在 \(k\) 上平展。则 \(X_\kappa\) 是由所有 \(x\in X\) 对应的 \(\operatorname{Spec}\kappa(x)\) 的不交并构成的形式 \(\kappa\)-簇。若 \(X'_\kappa\) 表示同一不交并给出的 \(\kappa\)-概形，则
\[
\begin{aligned}
  X_\kappa(\kappa_s)
  &=
  X'_\kappa(\kappa_s)
  \\
  &=
  \operatorname{Hom}_{\operatorname{Sch}/\kappa}
  \bigl(\operatorname{Spec}\kappa_s,X'_\kappa\bigr).
\end{aligned}
\]
令 \(\mathcal C\) 为平展 \(\kappa\)-概形构成的全子范畴
\[
  \operatorname{Sch}^{\mathrm{\acute{e}t}}/\kappa
  \subset
  \operatorname{Sch}/\kappa
\]
。已知
\[
  X'\longmapsto X'(\kappa_s)
\]
给出 \(\mathcal C\) 与 \(\Gamma\)-集合范畴 \(\mathcal C'\) 之间的等价\textup{（参见 SGA 1，V，§§7--8，或 \([DG70]\)，§I.4，6.4）}。
因而它诱导 \(\mathcal C\) 中的群对象范畴与 \(\mathcal C'\) 中的群对象范畴之间的等价。\(\mathcal C'\) 中的群对象，无非是带有 \(\Gamma\) 通过群自同构的连续作用的抽象群；我们称其为 \emph{\(\Gamma\)-群}。
结合 \SGARefLink{sga3:VIIB:1.6.E}{1.6.E} 中的等价
\[
  \operatorname{Vaf}^{\mathrm{\acute{e}t}}/k
  \xrightarrow{\ \sim\ }
  \operatorname{Vaf}^{\mathrm{\acute{e}t}}/\kappa
  \xrightarrow{\ \sim\ }
  \operatorname{Sch}^{\mathrm{\acute{e}t}}/\kappa,
\]
得到如下结论。
\medskip
\noindent\textbf{命题。}
\phantomsection
\label{sga3:VIIB:2.5:proposition}
{\itshape 设 \(k\) 是局部伪紧环，\(\kappa\) 是其剩余域，\(\kappa_s\) 是 \(\kappa\) 的可分闭包，且 \(\Gamma=\operatorname{Gal}(\kappa_s/\kappa)\)。
\begin{enumerate}
\item[\textup{(i)}] 函子 \(X\mapsto X(\kappa_s)\) 给出平展形式 \(k\)-簇范畴与 \(\Gamma\)-集合范畴之间的等价。
\item[\textup{(ii)}] 它诱导平展形式 \(k\)-群范畴与 \(\Gamma\)-群范畴之间的等价。
\end{enumerate}}
\medskip
\noindent\textbf{注 2.5.A。}
\phantomsection
\label{sga3:VIIB:2.5.A}
\footnote{\label{sga3:VIIB:nde:96}此注由编者添加。}
设 \(k\) 是域，\(k_s\) 是 \(k\) 的可分闭包，\(G\) 是平展形式 \(k\)-群，置
\[
  M=G(k_s).
\]
选取 \(\Gamma=\operatorname{Gal}(k_s/k)\) 在 \(M\) 中的轨道的一组代表 \(X\)。对每个 \(x\in X\)，令 \(\Gamma_x\) 为其稳定子，这是 \(\Gamma\) 的有限指数开子群，并置
\[
  L_x=k_s^{\Gamma_x}.
\]
则 \(L_x/k\) 的次数为 \([\Gamma:\Gamma_x]\)\textup{（例如参见 \([BAlg]\)，V，§10.10）}。
由上述范畴等价，\(G\) 的仿射代数 \(\mathcal A(G)\) 是各 \(L_x\) 的积，赋予积拓扑。因此，\(L_x\) 恰好是所有单商
\[
  \mathcal A(G)/\mathfrak m,
\]
其中 \(\mathfrak m\) 遍历 \(\mathcal A(G)\) 的开极大理想。由于这些商通过对偶与 \(\mathbf H(G)\) 的单子余代数对应，余代数 \(\mathbf H(G)\) 是点化的——即对每个单子余代数 \(C\) 有
\[
  \dim_kC=1
\]
——当且仅当对每个 \(x\) 有 \(L_x=k\)。此时 \(\mathcal A(G)\) 就是拓扑代数 \(k^M\)，故 \(\mathbf H(G)\) 是群代数 \(kM\)，且
\[
  M
  =
  \bigl\{
    x\in\mathbf H(G)
    \mathrel{\big|}
    \varepsilon(x)=1,\ \Delta(x)=x\otimes x
  \bigr\}
  =
  G(k).
\]
